\documentclass[10pt,a4paper]{amsart}
\usepackage[margin=19mm]{geometry}
\usepackage{amsmath,amssymb,mathtools}
\usepackage[scr=rsfs,cal=euler]{mathalfa}
\usepackage{xfrac}
\usepackage[unicode,hidelinks,bookmarks=false]{hyperref}
\newcommand{\ie}{i.e.\ }
\newcommand{\cf}{cf.\ }
\newcommand{\resp}{respectively\ }
\newcommand{\Subsection}{\subsection}
\providecommand{\nobreakdash}{\nobreak\hbox{-}}
\setlength{\emergencystretch}{2em}
\multlinegap0pt
\usepackage{mathtools}
\usepackage[shortlabels]{enumitem}
\usepackage{amssymb}
\usepackage{tikz}
\newtheorem{prop}{Proposition}
\newtheorem{thm}{Theorem}
\newcommand{\R}{\mathbb{R}}
\newcommand{\hs}{h^*}
\title{A very ample lattice polytope with a non-unimodal $h^*$-vector}
\author{Johannes Hofscheier, Vadym Kurylenko, Benjamin Nill}
\date{Source: arXiv:2608.21507v1 (CC BY 4.0)}
\begin{document}
\maketitle

\noindent\emph{Source and licence.} A very ample lattice polytope with a non-unimodal $h^*$-vector. Version arXiv:2608.21507v1 (CC BY 4.0), distributed by arXiv under the Creative Commons Attribution 4.0 International licence, http://creativecommons.org/licenses/by/4.0/, as recorded in the arXiv licence metadata for this exact version and shown on the abstract page https://arxiv.org/abs/2608.21507v1. Full licence notice: \texttt{licenses/arxiv-2608.21507-CC-BY-4.0.txt}.\par\smallskip

\noindent\textit{Editorial scope.} Counted unit: the article's thm thm:main (source0.tex lines 193--195). The article's complete original proof of it is delivered with it (source0.tex lines 197--209, one \texttt{proof} environment of 13 lines). No mathematics is written, repaired or re-derived: the statement and the proof are the article's own lines, in the article's own notation, and no retained line is substituted or re-flowed.  Retained uncounted as the source-local context the proof needs: lines 173--183.  Nothing was omitted from any retained span, so the package carries no omission record.  No retained line contains a citation, so the wrapper carries no bibliography and no bibliography entry.  No retained line contains a figure environment, an image-inclusion command or drawing code, so no image is delivered and the package declares no asset dependency.  Editorial formatting only, in addition: an AMS \texttt{amsart} preamble that restates the article's own declarations and macros used by the retained lines, namely the environments \texttt{prop}, \texttt{thm} on separate counters, together with the standard packages those lines need.  The retained environments print in this document as Proposition 1, Theorem 1 in order of appearance, whereas the article prints the same items with the numbering of its own full document.  The complete native source of the article as distributed in the arXiv e-print for this exact version is bundled unchanged as \texttt{sources/208254/source0.tex}.
\begin{prop}\label{prop}
    The lattice polytope $Q \subset \R^{31}$ has the following properties:
    \begin{enumerate}
        \item its dimension is $29$ and it has $60$ vertices,
        \item it is very ample but not IDP,
        \item its $h^*$-polynomial is
            \[
                h_Q^*(t)=1+30(t+\cdots+t^{14})+346t^{15}.
            \]
    \end{enumerate}
\end{prop}

\begin{thm} \label{thm:main}
    $Q \times Q$ is a $58$-dimensional very ample, non-IDP lattice polytope with $3600$ vertices and non-unimodal $h^*$-vector.
\end{thm}

\begin{proof}
    Clearly, Cartesian products preserve very ampleness as well as being non-IDP. Note that for all nonnegative integers $k$ we have $E_{Q \times Q}(k)=(E_Q(k))^2$. Therefore, for integers $r$ we get for the $r$th-coefficient of $h^*_{Q \times Q}(t)$
    \[h_r^*=\sum_{m=0}^r (-1)^{r-m} \binom{59}{r-m} E_Q(m)^2,\]
    where by Proposition~\ref{prop}
    \[E_Q(m)=\binom{m+29}{29}+\left(\sum_{i=1}^{14} 30 \binom{m+29-i}{29}\right) + 346 \binom{m+29-15}{29}.\] 
    Therefore, we can verify non-unimodality by evaluating
    
    \begin{align*}
        \hs_{25} &= 1631347108606059869973,\\
        \hs_{26} &= 1627940480253228812568,\\
        \hs_{27} &= 1628167830322084291128. \qedhere
    \end{align*}
\end{proof}

\end{document}
